\subsection{可优超的向量测度}

设 \(q\) 是 \(F\) 上的下半连续半范数。我们用 \(A_q'\) 表示由 \(z' \in F'\) （满足 \(|\langle z', x \rangle| \leqslant q(x)\) ，对一切 \(x \in F\) 者）所成之集。它是在 \(\mathbf{F}'\) 中由 \(\mathbf{x} \in \mathbf{F}\) （满足 \(q(\mathbf{x}) \leqslant 1\) 者）所成之集的极集；对每个 \(\mathbf{x} \in \mathbf{F}\) ， \(q(\mathbf{x}) = \sup_{\mathbf{z}' \in \hat{\mathrm{A}}_q'} |\langle \mathbf{x}, \mathbf{z}' \rangle|\) 。

定义 3. --- 设 m 是 T 上取值于 F 的向量测度。若 q 是 F 上的下半连续半范数，则称 m 是 q-可优超的，如果存在正测度 \(\mu\) ，使得 \(|z^{\prime} \circ m| \leqslant \mu\) （对每个 \(z^{\prime} \in A_{q}^{\prime}\) ）；这时把诸测度 \(|z^{\prime} \circ m|\) （ \(z^{\prime}\) 跑遍 \(A_{q}^{\prime}\) ）的上确界（Ch. III, §2, No. 4, 定理 3）记作 \(q(\mathbf{m})\) 。若 m 关于 F 上的每个连续半范数 q 都是 q-可优超的，则称 m 是可优超的。

若 \(\mathbf{m}\) 与 \(\mathbf{m}'\) 都是 \(q\)-可优超的，则立即可以看出 \(\mathbf{m} + \mathbf{m}'\) 也是 \(q\)-可优超的，并且

\[
q (\mathbf {m} + \mathbf {m} ^ {\prime}) \leqslant q (\mathbf {m}) + q (\mathbf {m} ^ {\prime}).
\]

当 F 是赋范空间、范数记作 \(|x|\) 时，说 m 可优超，就是说诸测度 \(|z'\circ m|\) （其中 \(|z'| \leqslant 1\) ）有同一个正测度作为上界\(^{1}\)；这时就把该测度族的上确界记作 \(|m|\) 。

若 \(\mathbf{F} = \mathbf{R}\) ，则测度 \(|m|\) （对应于欧几里得范数 \(|x|\) ， \(\mathbf{R}\) 上的）与 Ch. III, §1, No. 5 中记作 \(|m|\) 的测度一致。

命题 4. --- 设 \((\mathrm{F}_i)_{1 \leqslant i \leqslant n}\) 是有限族 Hausdorff 局部凸空间， \(\mathbf{F} = \prod_{i=1}^{n} \mathbf{F}_i\) 是它们的乘积， \(q_i\) ( \(1 \leqslant i \leqslant n\) ) 是 \(\mathbf{F}_i\) 上的下半连续半范数，而 \(q\) 是 \(\mathbf{F}\) 上由 \(q(\mathbf{x}_i, \ldots, \mathbf{x}_n) = \sum_{i=1}^{n} q_i(\mathbf{x}_i)\) 定义的半范数。若 \(\mathbf{m}_i\) ( \(1 \leqslant i \leqslant n\) ) 是 \(\mathbf{T}\) 上取值于 \(\mathbf{F}_i\) 的、 \(q_i\)-可优超的向量测度，则测度 \(\mathbf{m} = (\mathbf{m}_1, \ldots, \mathbf{m}_n)\) （取值于 \(\mathbf{F}\) ）是 \(q\)-可优超的。

因为，对偶 \(\mathbf{F}'\) 可以等同于 \(\prod_{i=1}^{n} \mathrm{F}_i'\) ，使得若 \(\mathbf{x} = (\mathbf{x}_i) \in \mathrm{F}\) ， \(\mathbf{z}' = (\mathbf{z}_i') \in \mathrm{F}'\) ，则 \(\langle \mathbf{x}, \mathbf{z}' \rangle = \sum_{i=1}^{n} \langle \mathbf{x}_i, \mathbf{z}_i' \rangle\) 。若 \(|\langle \mathbf{x}, \mathbf{z}' \rangle| \leqslant q(\mathbf{x})\) （对每个 \(\mathbf{x} \in \mathrm{F}\) ），则特别地 \(|\langle \mathbf{x}_i, \mathbf{z}_i' \rangle| \leqslant q_i(\mathbf{x}_i)\) （对 \(1 \leqslant i \leqslant n\) ），而其逆是显然的，这表明集 \(\mathrm{A}_q'\) 是诸 \(\mathrm{A}_{q_i}'\) 的乘积。由于按假设 \(|\mathbf{z}_i' \circ \mathbf{m}_i| \leqslant q_i(\mathbf{m}_i)\) （对 \(\mathbf{z}_i' \in \mathrm{A}_{q_i}'\) ），由此得出

\[
\left| \mathbf {z} ^ {\prime} \circ \mathbf {m} \right| \leqslant \sum_ {i = 1} ^ {n} \left| \mathbf {z} _ {i} ^ {\prime} \circ \mathbf {m} _ {i} \right| \leqslant \sum_ {i = 1} ^ {n} q _ {i} (\mathbf {m} _ {i})
\]

（对每个 \(\mathbf{z}' \in \mathrm{A}_q'\) ），这就证明了命题。

推论. --- 若空间 F 是有限维的，则每个取值于 F 的向量测度 m 都是可优超的。为使一个数值函数关于 m 本质可积，必须且只需它关于 \(|m|\) 本质可积（其中 \(|x|\) 表示 F 上的任一范数）。

命题 5. --- 设 q 是 F 上的下半连续半范数。设 m 是 q-可优超的测度， f 是关于 m 本质可积且使 \(\int f dm \in F\) 的函数。则

\[
q \left(\int f d \mathbf {m}\right) \leqslant \int^ {\bullet} | f | d q (\mathbf {m}).
\]

因为，

\[
q \left(\int f d \mathbf {m}\right) = \sup _ {\mathbf {z} ^ {\prime} \in \mathrm{A} _ {q} ^ {\prime}} \left| \left\langle \mathbf {z} ^ {\prime}, \int f d \mathbf {m} \right\rangle \right| \leqslant \sup _ {\mathbf {z} ^ {\prime} \in \mathrm{A} _ {q} ^ {\prime}} \int | f | d | \mathbf {z} ^ {\prime} \circ m | \leqslant \int^ {\bullet} | f | d q (\mathbf {m})
\]

这是根据 (1) 以及关系式 \(|\mathbf{z}' \circ \mathbf{m}| \leqslant q(\mathbf{m})\) （对 \(\mathbf{z}' \in \mathrm{A}_q'\) ）。

命题 6. --- 设 F 是拟完备的 Hausdorff 局部凸空间， m 是 T 上取值于 F 的可优超测度。若数值函数 f 关于所有测度 \(q(\mathbf{m})\) （其中 q 跑遍 F 上连续半范数所成之集）都本质可积，则 f 关于 m 本质可积，且 \(\int f dm \in F\) 。

我们将利用下面的辅助结果。设 \((\mu_{\iota})_{\iota\in\mathrm{I}}\) 是 T 上正测度的递增有向族。用 \(\mathcal{L}^{1}\big((\mu_{\iota})_{\iota\in\mathrm{I}}\big)\) 表示 T 上有限数值函数中本质 \(\mu_{\iota}\)-可积（对每个 \(\iota\in I\) ）者所成的向量空间，赋以由半范数 \(f\mapsto\mu_{\iota}(|f|)\;\ (\iota\in\mathrm{I})\) 定义的拓扑。设 \(L_{0}\) 是 \(\mathcal{L}^{1}\big((\mu_{\iota})_{\iota\in\mathrm{I}}\big)\) 中由诸乘积 \(g\varphi_{K}\) 生成的线性子空间，其中 g 跑遍 T 上连续有限数值函数所成之集， K 跑遍 T 的紧子集所成之集。

引理 1. --- 当 \(\mathcal{L}_0\) 与 \(\mathcal{K}(\mathrm{T})\) 赋以 \(\mathcal{L}^1\big((\mu_\iota)_{\iota \in \mathrm{I}}\big)\) 的拓扑诱导的拓扑时：

\begin{enumerate}
\def\labelenumi{\alph{enumi})}
\item
  \(\mathcal{L}_0\) 的每个元素都属于 \(\mathcal{K}(\mathrm{T})\) 的某个有界子集的闭包；
\item
  \(\mathcal{L}^1\big((\mu_\iota)_{\iota \in \mathrm{I}}\big)\) 的每个元素都属于 \(\mathcal{L}_0\) 的某个有界子集的闭包。
\end{enumerate}

为证 a)，可以只考虑形如 \(f = g\varphi_{\mathrm{K}}\) 的元素（ \(g \in \mathcal{C}(\mathrm{T})\) ， K 为 T 中的紧集）。立即看出（借助 Urysohn 定理）， \(f\) 属于由形如 \(gh\) 的函数组成的集 B 的闭包，其中 \(h\) 跑遍 T 到 {[}0, 1{]} 中的连续映射、在 K 上等于 1 且在 K 的一个固定紧邻域 H 的余集上等于 0 者。此外，集 B 是有界的，因为 \(\mu_{\iota}(|gh|) \leqslant \mu_{\iota}(|f\varphi_{\mathrm{H}}|)\) （对每个具有前述性质的函数 h）。

现在来证 b)；可以只考虑形如 \(f \geqslant 0\) 的元素（ \(\mathcal{L}^1((\mu_\iota)_{\iota \in \mathrm{I}})\) 中的）。对每个 \(\iota \in \mathrm{I}\) 与每个 \(\varepsilon > 0\) ，存在一个紧子集 \(\mathrm{K}(\iota, \varepsilon)\) （ \(\mathrm{T}\) 的），使得 \(f\) 在 \(\mathrm{K}(\iota, \varepsilon)\) 上的限制连续且 \(|\mu_\iota(|f - f\varphi_{\mathrm{K}(\iota, \varepsilon)}|) \leqslant \varepsilon\) 。显然 \(f\) 属于集 \(\mathrm{C}\) ——由诸 \(f\varphi_{\mathrm{K}(\iota, \varepsilon)}\) （其中 \(\iota \in \mathrm{I}, \varepsilon > 0\) ）组成——的闭包。借助 Urysohn 定理，集 \(\mathrm{C}\) 含于 \(\mathcal{L}_0\) ；此外它是有界的，因为 \(\mu_\varkappa(f\varphi_{\mathrm{K}(\iota, \varepsilon)}) \leqslant \mu_\varkappa(f)\) （对所有 \(\iota \in \mathrm{I}, \varkappa \in \mathrm{I}\) 与 \(\varepsilon > 0\) ）。

现在来证命题 6：对每个函数 \(g \in \mathcal{K}(\mathrm{T})\) 与每个连续半范数 \(q\) （ \(\mathrm{F}\) 上的），有 \(q\big(\int g dm\big) \leqslant \int |g| d\big(q(\mathbf{m})\big)\) （命题 5），这蕴含映射 \(g \mapsto \int g dm\) （ \(\mathcal{K}(\mathrm{T})\) 到 \(\mathrm{F}\) 中的）当 \(\mathcal{K}(\mathrm{T})\) 赋以 \(\mathcal{L}^1\big((q(\mathbf{m}))_{q \in \mathbb{Q}}\big)\) 的拓扑诱导的拓扑时（Q 为 \(\mathrm{F}\) 上连续半范数之集）是连续的。因此，由前述引理与 TVS, III, §1, No. 6 的命题 10，这一映射可以连续地延拓，先到 \(v_0\) —— \(\mathcal{L}_0\) 到 \(\mathrm{F}\) 中的连续线性映射——再到 \(v\) —— \(\mathcal{L}^1\big((q(\mathbf{m}))_{q \in \mathbb{Q}}\big)\) 到 \(\mathrm{F}\) 中的连续线性映射。此外，对每个 \(\mathbf{z}' \in \mathrm{F}'\) ，关系式 \(\langle \mathbf{z}', v(f) \rangle = \int f d(\mathbf{z}' \circ \mathbf{m})\) （按 \(v\) 的定义）对每个 \(f \in \mathcal{K}(\mathrm{T})\) 成立；由于 \(|\mathbf{z}' \circ \mathbf{m}| \leqslant q(\mathbf{m})\) （对 \(q(\mathbf{z}) = |\langle \mathbf{z}', \mathbf{z} \rangle|\) 而言），映射 \(f \mapsto \int f d(\mathbf{z}' \circ \mathbf{m})\) 在 \(\mathcal{L}^1\big((q(\mathbf{m}))_{q \in \mathbb{Q}}\big)\) 上连续，因此再由连续性，关系式 \(\langle \mathbf{z}', v(f) = \int f d(\mathbf{z}' \circ \mathbf{m})\) 对每个函数 \(f \in \mathcal{L}^1\big((q(\mathbf{m}))_{q \in \mathbb{Q}}\big)\) 成立。由此得出 \(\nu(f) = \int f dm\) ，证明完成。

